\section{Overview and interfaces}
\label{sec:overview}

\begingroup\small

\subsection{Committed and virtual witness views}
\label{sec:public-objects}

The PCS receives a committed binary witness $\bm f$.  Before proving, the
developer fixes the matrix layouts used by the commitment and the PIOP.
The same configuration may optionally include a public virtualization map:

\begin{itemize}[leftmargin=1.4em]
 \item \textbf{Direct.}
       The virtualization map is omitted and the PIOP uses the
       committed witness directly, so $\bm h=\bm f$.
 \item \textbf{Virtualized.}
       The PIOP uses a binary witness $\bm h$ related to the
       committed witness by a nontrivial developer-selected public matrix
       $\bm M$.
\end{itemize}

Throughout, $\bm f$ is the commitment-facing vector and $\bm h$ is the
PIOP-facing vector.  When $\bm M$ is omitted, every uniform formula
below interprets it as the identity map.

The virtual witness is
\begin{equation}
 \boxed{\quad
 \pi_2(\bm h)=\bm M\pi_2(\bm f).
 \quad}
 \label{eq:virtual-witness}
\end{equation}
Here $\pi_2:\mathbb Z\to\Ftwo$, $x\mapsto x\bmod 2$, acts coordinatewise.
Since $\bm f$ and $\bm h$ are binary, Equation~\eqref{eq:virtual-witness}
uniquely determines $\bm h$.  If virtualization is enabled, the developer
fixes $\bm M$ before proving; it is not chosen by the prover.

Only $\bm f$ is committed.  The developer-configured matrices $F$ and $H$ are
reshapings of $\bm f$ and $\bm h$, not additional witnesses.  The commitment is
\begin{equation}
 \bm f\in\bitsset^{n_f},\qquad
 \operatorname{vec}(F):=\pi_2(\bm f),\qquad
 P:=\Pack(F),\qquad
 O:=\Enc_{\mathcal C}(P),\qquad
 \rho:=\Root(O).
 \label{eq:overview-commitment}
\end{equation}
Neither $H$, the layout of $\bm h$, nor any value derived from $H$ is encoded
in $O$.

\paragraph{Adjoint convention.}
For the standard inner product over a field $\mathbb E$, the adjoint of a
matrix $\bm T$ is its transpose $\bm T^{\mathsf T}$:
\begin{equation}
 \langle\bm u,\bm T\bm v\rangle_{\mathbb E}
 =\langle\bm T^{\mathsf T}\bm u,\bm v\rangle_{\mathbb E}.
 \label{eq:adjoint-convention}
\end{equation}
Throughout this specification, applying the adjoint means using this identity
to move a public matrix across the inner product.  It is a deterministic
rewrite and sends no protocol message.

\subsection{PIOP linear-claim handoff}
\label{sec:piop-handoff}

The end-to-end proof begins with a constraint-satisfaction relation over $\Fq$
evaluated on the virtual witness $\pi_q(\bm h)$.  The PIOP reduces that relation
to a general Lin-BitsDual claim consumed by F2Z.  The input Sumcheck phase of
$\Pi_1$ shapes that claim before the integer-fold phase:
\begin{equation}
 \begin{aligned}
  (x;\pi_q(\bm h))\in\mathcal R_{\rm CS}
  &\xRightarrow{\ \Pi_{\rm PIOP}\ }
    \left\langle\bm v,\pi_q(\bm h)\right\rangle_{\Fq}=\mu \\
  &\xRightarrow{\ \Pi_1\ }
    \mathcal R_{\rm PESAT}[\bm h] \\
  &\xRightarrow{\ \mathrm{F2Z}\ }
    \text{opening against $\rho$}.
 \end{aligned}
 \label{eq:piop-f2z-flow}
\end{equation}
The PIOP is defined by the constraint system and is not part of F2Z itself.
The end-to-end interaction includes it as the subprotocol that produces the
F2Z handoff.  F2Z receives the arbitrary public coefficient $\bm v$ and does
not inspect the equations or reductions that produced it.  The checked input
Sumcheck is internal to $\Pi_1$; it does
not factor $\bm v$, but replaces the original scalar claim by an
evaluation-shaped intermediate claim before producing F2Z PESAT instances.

\endgroup

\section{Instance and parameters}
\label{sec:instance}

\begingroup\small

Section~\ref{sec:instance} is the canonical registry for persistent and
cross-phase F2Z notation.  A definition gives the exact value denoted by a
symbol, its type gives the ambient set or interface class, and its rule records
the normative constraint on its use.  Section~\ref{sec:encoding} specifies how
the transcript derives the non-input symbols registered here.

\newcommand{\DefTableHead}{%
  \toprule
  \cellcolor{specnavy}\textcolor{white}{\textbf{Symbol}} &
  \cellcolor{specblue}\textcolor{white}{\textbf{Definition}} &
  \cellcolor{spectableteal}\textcolor{white}{\textbf{Type}} &
  \cellcolor{specnavy}\textcolor{white}{\textbf{Rule}} \\
  \midrule}

\newcommand{\DefTableSetup}{%
  \setlength{\tabcolsep}{3pt}%
  \renewcommand{\arraystretch}{1.10}%
  \arrayrulecolor{spectableline}%
  \rowcolors{2}{white}{spectablerow}}

\subsection{Algebraic domains and canonical maps}
\label{sec:domains}

\begingroup
\scriptsize
\DefTableSetup
\begin{tabularx}{\linewidth}{@{}>{\color{specnavy}\raggedright\arraybackslash}p{0.78in}
 >{\raggedright\arraybackslash}X
 >{\color{spectableteal}\raggedright\arraybackslash}p{1.25in}
 >{\raggedright\arraybackslash}X@{}}
\DefTableHead
$\mathbb Z$ & Exact integer domain & Ring \john{don't need} &\\
\midrule
$\Ftwo$ & Binary field & $\mathbb F_2$ &\\
\midrule
$\Fq$ & $q=2^{100}-15$ & Finite field &\\
\midrule
$\K$ & $\Ftwo[X]/(X^{128}+X^7+X^2+X+1)$ & $\mathbb F_{2^{128}}$ & Used for packing, products, GKR, and openings. \\
\midrule
$\pi_2$ & $x\mapsto x\bmod2$ & $\mathbb Z\to\Ftwo$ & \\
\midrule
$\operatorname{lift}_2(b)$ & The unique $u\in\{0,1\}\subset\mathbb Z$ satisfying $\pi_2(u)=b$ & $\Ftwo\to\{0,1\}$ & Canonical bit representative; not an inverse of $\pi_2$ on all of $\mathbb Z$. \\
\midrule
$\pi_q$ & $x\mapsto x\bmod q$ & $\mathbb Z\to\Fq$ & \\
\midrule
$\operatorname{lift}_q$ & $\operatorname{lift}_q(a)$ is the unique integer in $[0,q)$ satisfying $\pi_q(\operatorname{lift}_q(a))=a$ & $\Fq\to[0,q)$ & Canonical integer lift; no alternative representative is accepted. \\
\midrule
$\iota_{2\to\K}$ & Canonical subfield embedding & $\Ftwo\hookrightarrow\K$ & Applied coordinatewise when a binary witness enters a $\K$-claim. \\
\midrule
$\beta_v$ & $X^v$ in the fixed polynomial basis of $\K$ & Element of $\K$ & $0\le v<128$, ordered by increasing $v$. \\
\bottomrule
\end{tabularx}
\endgroup

\subsection{Dimensions and index conventions}
\label{sec:dimensions}

\begingroup
\scriptsize
\DefTableSetup
\begin{tabularx}{\linewidth}{@{}>{\color{specnavy}\raggedright\arraybackslash}p{0.78in}
 >{\raggedright\arraybackslash}X
 >{\color{spectableteal}\raggedright\arraybackslash}p{1.25in}
 >{\raggedright\arraybackslash}X@{}}
\DefTableHead
$n_f$ & Length of the committed witness & $\mathbb N$ & $n_f=2^{m_f}$ and $128\mid n_f$. \\
\midrule
$n_h$ & Length of the virtual witness & $\mathbb N$ & $n_h=2^{m_h}$. \\
\midrule
$m_f$ & Boolean dimension of $\bm f$ & $\mathbb N$ & \\
\midrule
$m_h$ & Boolean dimension of $\bm h$ & $\mathbb N$ & \\
\midrule
$d_{F,r},d_{F,c}$ & Row and column dimensions of $F$ & $\mathbb N^2$ & $d_{F,r}+d_{F,c}=m_f$. \\
\midrule
$d_r,d_c$ & Row and column dimensions of $H$ & $\mathbb N^2$ & $d_r+d_c=m_h$ and $2^{d_r}q<2^{127}$. \\
\midrule
$m_p$ & Packed-message dimension & $\mathbb N$ & $m_p=m_f-7$. \\
\midrule
$r$ & Row index of $H$ & $0\le r<2^{d_r}$ &  \\
\midrule
$c$ & Column index of $H$ & $0\le c<2^{d_c}$ & \\
\midrule
$j$ & Packed-element index & $0\le j<n_f/128$ & \\
\midrule
$v$ & Basis-coordinate index in $K$ & $0\le v<128$ & index to reference a coefficient for a packed galios number \\
\bottomrule
\end{tabularx}
\endgroup

\subsection{Witnesses, layouts, and virtualization}
\label{sec:witness-layouts}

\begingroup
\scriptsize
\DefTableSetup
\begin{tabularx}{\linewidth}{@{}>{\color{specnavy}\raggedright\arraybackslash}p{0.78in}
 >{\raggedright\arraybackslash}X
 >{\color{spectableteal}\raggedright\arraybackslash}p{1.25in}
 >{\raggedright\arraybackslash}X@{}}
\DefTableHead
$\bm f$ & Commitment-facing binary witness & $\bitsset^{n_f}$ & The PCS commits only to this vector. \\
\midrule
$F$ & Developer-configured reshaping of $\pi_2(\bm f)$ & $\Ftwo^{2^{d_{F,r}}\times2^{d_{F,c}}}$ & $\operatorname{vec}(F)=\pi_2(\bm f)$. The developer reshapes it in terms of $2^{d_{F,r}}\times2^{d_{F,c}}$. \\
\midrule
$\bm h$ & PIOP-facing binary witness & $\bitsset^{n_h}$ & Equals $\bm f$ when virtualization is omitted. \\
\midrule
$H$ & Developer-configured reshaping of $\pi_2(\bm h)$ & $\Ftwo^{2^{d_r}\times2^{d_c}}$ & $\operatorname{vec}(H)=\pi_2(\bm h)$. \\
\midrule
$H_{\mathbb Z}$ & Entrywise $\operatorname{lift}_2(H)$ & $\{0,1\}^{2^{d_r}\times2^{d_c}}\subset\mathbb Z^{2^{d_r}\times2^{d_c}}$ & Used in exact integer folds. \\
\midrule
$H_q$ & Entrywise $\pi_q(H_{\mathbb Z})$ & $\Fq^{2^{d_r}\times2^{d_c}}$ & Used in the Lin-BitsDual claim and the input Sumcheck. \\
\midrule
$H_{\K}$ & Entrywise $\iota_{2\to\K}(H)$ & $\K^{2^{d_r}\times2^{d_c}}$ & Used in PESAT, GKR, and MLE claims. \\
\midrule
$\operatorname{vec}$ & Flattening of a configured matrix layout & Matrix-to-vector map & The developer fixes its coordinate order before proving. \\
\midrule
$\bm M$ & Effective virtualization map & $\Ftwo^{n_h\times n_f}$ & Supplied by the developer when virtualization is enabled.  Otherwise $n_h=n_f$, both layouts preserve the same flattened coordinate order, $\bm h=\bm f$, and $\bm M:=I$. \\
\midrule
$\bm M_{\K}$ & Scalar extension of the effective $\bm M$ & $\K^{n_h\times n_f}$ & Its transpose is the adjoint from Equation~\eqref{eq:adjoint-convention}. \\
\midrule
$\bm f_{\K},\bm h_{\K}$ & Coordinatewise images under $\iota_{2\to\K}$ & $\K^{n_f}$ and $\K^{n_h}$ & Used only in $\K$-linear claims. \\
\bottomrule
\end{tabularx}
\endgroup

The matrices $F$ and $H$ are views of the witness vectors, not additional
witnesses.  Their shapes and coordinate orders are developer configuration
choices fixed before proving.  The configured Boolean address of $H$ is
$z=(r,c)\in\bitsset^{d_r}\times\bitsset^{d_c}$; its vectorization and every
multilinear extension in this specification use this same row/column bit
split.  In particular, reshaping a vector $\bm u\in\Fq^{n_h}$ into a matrix
$U$ means $\operatorname{vec}(U)=\bm u$ under that configured order, not an
independently chosen row-major convention.

\subsection{Packing and commitment}
\label{sec:packing-commitment}

\begingroup
\scriptsize
\DefTableSetup
\begin{tabularx}{\linewidth}{@{}>{\color{specnavy}\raggedright\arraybackslash}p{0.78in}
 >{\raggedright\arraybackslash}X
 >{\color{spectableteal}\raggedright\arraybackslash}p{1.25in}
 >{\raggedright\arraybackslash}X@{}}
\DefTableHead
$\Pack$ & Packs consecutive groups of $128$ coordinates into $\K$ & $\Ftwo^{n_f}\to\K^{n_f/128}$ & Uses the ordered basis $(\beta_v)_{v<128}$. \\
\midrule
$P$ & $P[j]=\sum_{v=0}^{127}\operatorname{vec}(F)[128j+v]\beta_v$ & $\K^{n_f/128}$ & Derived only from the configured view $F$. \\
\midrule
$F_v$ & $F_v[j]=\operatorname{vec}(F)[128j+v]$ & $\Ftwo^{2^{m_p}}$ & The $v$-th bit slice of the packed message, so $P[j]=\sum_vF_v[j]\beta_v$. \\
\midrule
$\mathcal C$ & Configured encoding code & Public code &  \\
\midrule
$\Enc_{\mathcal C}$ & Canonical encoder for $\mathcal C$ & $\K^{n_f/128}\to\K^{n_O}$ & \\
\midrule
$O$ & $\Enc_{\mathcal C}(P)$ & $\K^{n_O}$ & Encodes $P$ and never the virtual view $H$. \\
\midrule
$\Root$ & Configured Merkle-root algorithm & $\K^{n_O}\to\bitsset^{\lambda}$ & Constructs a Merkle tree on $\Enc_{\mathcal C}(P)$. \\
\midrule
$\rho$ & $\Root(O)$ & $\bitsset^{\lambda}$ & Public commitment to $O$. \\
\bottomrule
\end{tabularx}
\endgroup

\subsection{Structured ring-switch coefficients}
\label{sec:ring-switch-coefficients}

\begingroup
\scriptsize
\DefTableSetup
\begin{tabularx}{\linewidth}{@{}>{\color{specnavy}\raggedright\arraybackslash}p{0.78in}
 >{\raggedright\arraybackslash}X
 >{\color{spectableteal}\raggedright\arraybackslash}p{1.25in}
 >{\raggedright\arraybackslash}X@{}}
\DefTableHead
$A_f[v,j]$ & $\bm a_f[128j+v]$ in the configured commitment order & $\K^{128\times2^{m_p}}$ & Pack-boundary view of the public coefficient produced by the adjoint. \\
\midrule
$\Theta$ & $\{0,\ldots,t_{\rm sw}-1\}$ in canonical order & Finite nonempty index set & Derived by the developer configuration from $\bm a_f$, with $t_{\rm sw}<2^{32}$. \\
\midrule
$A_\theta,B_\theta$ & $A_\theta[v]\in\K$ and $B_\theta[j]\in\K$ & $\K^{128}$ and $\K^{2^{m_p}}$ & Must satisfy $A_f[v,j]=\sum_{\theta\in\Theta}A_\theta[v]B_\theta[j]$ for every $(v,j)$. \\
\midrule
$\mathsf{RingDecomp}$ & $\bm a_f\mapsto(\Theta,(A_\theta,B_\theta)_{\theta\in\Theta})$ & Public deterministic map & The verifier derives the same ordered decomposition; no proof field is used. \\
\bottomrule
\end{tabularx}
\endgroup

Every coefficient admits the canonical fallback
$\Theta=\{0,\ldots,127\}$,
$A_\theta[v]=1$ when $v=\theta$ and $0$ otherwise, and
$B_\theta[j]=A_f[\theta,j]$.  A developer configuration may instead provide a
smaller structured decomposition, such as the translated-equality terms
induced by rotations, shifts, or other public witness wiring.

\subsection{Lin-BitsDual public claim}
\label{sec:linbits-instance}

\begingroup
\scriptsize
\DefTableSetup
\begin{tabularx}{\linewidth}{@{}>{\color{specnavy}\raggedright\arraybackslash}p{0.78in}
 >{\raggedright\arraybackslash}X
 >{\color{spectableteal}\raggedright\arraybackslash}p{1.25in}
 >{\raggedright\arraybackslash}X@{}}
\DefTableHead
$\bm v$ & Public coefficient vector & $\Fq^{n_h}$ & Uses the configured coordinate order of $H$. \\
\midrule
$V$ & $\operatorname{vec}(V)=\bm v$ in the configured order of $H$ & $\Fq^{2^{d_r}\times2^{d_c}}$ & Arbitrary; no rank or separability condition is imposed. \\
\midrule
$\mu$ & Claimed inner product $\langle\bm v,\pi_q(\bm h)\rangle_{\Fq}=\langle\bm v,\operatorname{vec}(H_q)\rangle_{\Fq}$ & $\Fq$ & Initial target of the input Sumcheck. \\
\midrule
$\mathcal I_{\rm LinBitsDual}$ & $(\bm v,\mu)$ together with the developer configuration & Public F2Z instance & Binds the arbitrary input claim, both layouts, and the effective virtualization map. \\
\midrule
$\bm s_h=(\bm s_r,\bm s_c)$ & Input-Sumcheck terminal point & $\Fq^{m_h}$ & Transcript-derived with $|\bm s_r|=d_r$ and $|\bm s_c|=d_c$. \\
\midrule
$\kappa$ & $\widetilde{\bm v}(\bm s_h)$ & $\Fq$ & Computed by the verifier; never supplied by the prover. \\
\midrule
$\tau$ & Input-Sumcheck terminal value & $\Fq$ & Reduced target satisfying $\kappa\widetilde H_q(\bm s_h)=\tau$. \\
\bottomrule
\end{tabularx}
\endgroup

After the input Sumcheck, the verifier locally derives
$a_r:=\eq_{\Fq}(r,\bm s_r)$ and
$a'_c:=\kappa\eq_{\Fq}(c,\bm s_c)$.  Consequently,
\begin{equation}
 W[r,c]:=\kappa\eq_{\Fq}((r,c),\bm s_h)
        =\eq_{\Fq}(r,\bm s_r)\,
         \kappa\eq_{\Fq}(c,\bm s_c)=a_ra'_c,
 \label{eq:derived-core-coefficient}
\end{equation}
only the new coefficient $W$ is separable.  Equation~\eqref{eq:derived-core-coefficient}
is a consequence of the equality polynomial at the Sumcheck terminal point;
it is not a decomposition of the arbitrary input matrix $V$.

\subsection{Integer lifts and fold parameters}
\label{sec:core-parameters}

\begingroup
\scriptsize
\DefTableSetup
\begin{tabularx}{\linewidth}{@{}>{\color{specnavy}\raggedright\arraybackslash}p{0.86in}
 >{\raggedright\arraybackslash}X
 >{\color{spectableteal}\raggedright\arraybackslash}p{1.27in}
 >{\raggedright\arraybackslash}X@{}}
\DefTableHead
$\gamma_r$ & $\operatorname{lift}_q(a_r)$ & $[0,q)\subset\mathbb Z$ & Unique canonical integer lift of $a_r$. \\
\midrule
$\widehat\mu_c$ & $\sum_{r<2^{d_r}}\gamma_rH_{\mathbb Z}[r,c]$ & $[0,2^{127})\subset\mathbb Z$ & The bound follows from $2^{d_r}q<2^{127}$ and is required for exponent injectivity. \\
\midrule
$\mu_c$ & Prover-supplied claimed fold & $[0,2^{127})\subset\mathbb Z$ & Accepted only if its exponent matches PESAT and $\tau=\sum_ca'_c\pi_q(\mu_c)$. \\
\midrule
$g$ & Product generator & $\K^\times$ & $\operatorname{ord}(g)=2^{128}-1$, making exponentiation injective on accepted folds. \\
\midrule
$\nu_c$ & $g^{\mu_c}$ & $\K^\times$ & Must equal $\prod_r\bigl(1+(g^{\gamma_r}-1)H_{\K}[r,c]\bigr)$. \\
\bottomrule
\end{tabularx}
\endgroup

\subsection{Protocol interfaces}
\label{sec:protocol-interfaces}

\begingroup
\scriptsize
\DefTableSetup
\begin{tabularx}{\linewidth}{@{}>{\color{specnavy}\raggedright\arraybackslash}p{0.9in}
 >{\raggedright\arraybackslash}X
 >{\raggedright\arraybackslash}X@{}}
\toprule
\cellcolor{specnavy}\textcolor{white}{\textbf{Symbol}} &
\cellcolor{specblue}\textcolor{white}{\textbf{Definition}} &
\cellcolor{specnavy}\textcolor{white}{\textbf{Rule}} \\
\midrule
$\Pi_1$ & $\begin{gathered}\langle\bm v,\pi_q(\bm h)\rangle_{\Fq}=\mu\\[-2pt] \Longrightarrow\quad g^{\mu_c}=\displaystyle\prod_r\left(1+(g^{\gamma_r}-1)H_{\K}[r,c]\right)\end{gathered}$ & Runs the input Sumcheck, derives the canonical integer weights, and checks the fold bounds, exponentiation, and reconstruction of the reduced scalar claim. \\
\midrule
$\Pi_{\rm GKR}$ & $\begin{gathered}g^{\mu_c}=\displaystyle\prod_r\left(1+(g^{\gamma_r}-1)H_{\K}[r,c]\right)\\[-2pt] \Longrightarrow\ \displaystyle\sum_z\omega_{\xi,z}H_{\K}[z]=\sigma\end{gathered}$ & \\
\midrule
$\Pi_{\rm MLE}$ & $\begin{gathered}\displaystyle\sum_z\omega_{\xi,z}H_{\K}[z]=\sigma\\[-2pt] \Longrightarrow\ \widetilde H_{\K}(\bm r_h)=\eta\end{gathered}$ & Typically reduced by $\mathsf{LinCheck}_H$. \\
\midrule
$\operatorname{Adj}_{\bm M}$ & $\langle\bm u,\bm M_{\K}\bm v\rangle_{\K}\longmapsto\langle\bm M_{\K}^{\mathsf T}\bm u,\bm v\rangle_{\K}$ & \\
\midrule
$\Pi_{\rm iopp}$ & $\begin{gathered}\langle\bm a_f,\bm f_{\K}\rangle_{\K}=\eta\\[-2pt] \Longrightarrow\ \langle B,P\rangle_{\K}=\beta\ \Longrightarrow\ \mathsf{accept}\end{gathered}$ & Performs one ring switch and one recursive Ligerito opening. \\
\bottomrule
\end{tabularx}
\endgroup

\endgroup
